% =============================================================================
%  RM75 可微逆可达算子
%  IRDPLAYGROUND 设计的完整数学说明（中文版）
% =============================================================================
% Overleaf / 本机推荐：
%   Compiler = XeLaTeX
%   xelatex ird_operator_mathematics_zh.tex
% 使用 ctex + Fandol（TeX Live / Overleaf 自带），不依赖 Noto CJK。
% 英文对照：ird_operator_mathematics.tex（pdflatex，无 fontspec）
% =============================================================================
\documentclass[UTF8,a4paper,11pt,fontset=fandol]{ctexart}

\usepackage[margin=2.4cm]{geometry}
\usepackage{amsmath,amssymb,amsthm,mathtools}
\usepackage{bm}
\usepackage{booktabs}
\usepackage{enumitem}
\usepackage{hyperref}
\usepackage{xcolor}
\usepackage{listings}
\usepackage{graphicx}
\usepackage{cleveref}

\hypersetup{
  colorlinks=true,
  linkcolor=blue!50!black,
  citecolor=blue!50!black,
  urlcolor=blue!60!black
}

\newtheorem{definition}{定义}
\newtheorem{remark}{注记}
\newtheorem{proposition}{命题}

\newcommand{\SE}{\mathrm{SE}}
\newcommand{\SO}{\mathrm{SO}}
\newcommand{\R}{\mathbb{R}}
\newcommand{\Tcp}{T_{\mathrm{tcp}}}
\newcommand{\Taxis}{T_{\mathrm{axis}}}
\newcommand{\Tfl}{T_{\mathrm{flange}}}
\newcommand{\softmin}{\operatorname{softmin}}
\newcommand{\softmax}{\operatorname{softmax}}
\newcommand{\CVaR}{\mathrm{CVaR}}
\newcommand{\LogSumExp}{\operatorname{LogSumExp}}

\title{RM75 的可微逆可达性：\\
IRDPLAYGROUND 算子的完整数学设计}
\author{IRDPLAYGROUND 技术说明（中文版）}
\date{}

\begin{document}
\maketitle

\begin{abstract}
本文给出 \texttt{ird\_playground} 中逆可达（IRD）算子的自洽数学说明。
该算子回答的问题比经典离散 IRD 表更窄：在给定医学上合法的工具中心点（TCP）位姿、
以及由导轨诱导的机器人基座（J1 轴）坐标系时，哪一个光滑方向能提高经过碰撞校验的
机械臂可达裕度（clearance），并可选择地在局部位置/姿态区域内成立？
我们将设计与四篇基础文献相对应
（Vahrenkamp 等 2013；Rudorfer 2024/RM4D；Murooka 等 2025；Chiu 1988），
推导代码中使用的每一个聚合公式，并解释为何自动微分能给出可用的 IRD 梯度、
为何对局部姿态锥的条件查询良定义，以及为何同一场无需重训练即可支持轨迹参数更新。
\end{abstract}

\tableofcontents
\newpage

% -----------------------------------------------------------------------------
\section{问题陈述与经典 IRD}
\label{sec:problem}

\subsection{可达性与逆可达性}
\label{sec:classical}

设机器人构型为 \(q\in\mathcal{Q}\subset\R^{n}\)，正运动学将构型映到 TCP 位姿：
\begin{equation}
\label{eq:fk}
  \mathrm{FK}:\mathcal{Q}\to\SE(3),\qquad
  q\mapsto \Tcp(q).
\end{equation}
其中 \(\SE(3)\) 是刚体位姿的特殊欧氏群
（旋转 \(R\in\SO(3)\) 与平移 \(p\in\R^{3}\)）。

\begin{definition}[二值可达性]
位姿 \(r\in\SE(3)\) 称为\emph{可达}，若存在无碰撞构型 \(q\)，使得
\(\mathrm{FK}(q)=r\)（在 IK 容差内）且满足关节限位。记
\begin{equation}
\label{eq:binary-reach}
  \mathcal{R}
  \;=\;
  \bigl\{\, r\in\SE(3)
  \;\big|\;
  \exists\, q\in\mathcal{Q}_{\mathrm{free}}:
  \mathrm{FK}(q)=r
  \,\bigr\}.
\end{equation}
\end{definition}

Vahrenkamp 等（2013）的经典\emph{可达分布}（RD）在相对固定机器人基座的工作空间
体素化中存储质量分数，例如映射到该体素的关节构型测度。
\emph{逆可达分布}（IRD）对同一数据取逆：给定固定 TCP（例如抓取），
得到使该 TCP 可达的候选基座位姿分布。

形式地，若 \(t_{i}=(\Tcp)_{i}^{-1}\) 表示 RD 中存储的基座到 TCP 样本，
则对世界系 TCP \(p\in\SE(3)\)，对应候选基座位姿为
\begin{equation}
\label{eq:ird-inversion}
  b_{i}(p) \;=\; p\, t_{i}^{-1}
  \;\in\;\SE(3)
\end{equation}
（平面基座放置时可再投影到 \(\SE(2)\)）。
\(b_{i}(p)\) 的 IRD 分数继承自 \(t_{i}\) 的 RD 质量。
这是离线离散查表。

\subsection{IRDPLAYGROUND 实际计算的内容}
\label{sec:our-question}

我们的一阶段算子\emph{不}替代多种子 IK 或最终碰撞校验。它回答：

\begin{quote}
对医学上合法的 TCP 位姿 \(\Tcp\) 与导轨诱导的 J1 轴坐标系 \(\Taxis\)，
哪一个光滑方向能提高经过碰撞校验的机械臂可达 clearance，
并在局部位置/姿态不确定或任务区域内仍然成立？
\end{quote}

记感兴趣的决策变量为 \(x\)（导轨坐标、血管表面 TCP 转角、轨迹控制点，\ldots）。
我们构造标量
\begin{equation}
\label{eq:operator-goal}
  C(x)
  \;=\;
  \mathcal{A}\bigl[\, f\bigl(c(\Tcp(x),\Taxis(x))\bigr)\,\bigr]
\end{equation}
其中：
\begin{itemize}[leftmargin=1.4em]
  \item \(c(\cdot)\) 是光滑的 RM4D 风格图卡（第~\ref{sec:chart}~节）；
  \item \(f\) 是学习得到的有符号 clearance 场（第~\ref{sec:field}~节）；
  \item \(\mathcal{A}\) 是可微集合聚合器
        （Region~A、TaskCone、SetQuery 或 Trajectory；
        第~\ref{sec:region}--\ref{sec:traj}~节）。
\end{itemize}
基于梯度的更新使用贯穿整条 Torch 计算图的自动微分 \(\nabla_{x} C\)。

% -----------------------------------------------------------------------------
\section{文献对应：各篇贡献了什么}
\label{sec:papers}

\begin{table}[h]
\centering
\begin{tabular}{@{}llp{6.6cm}@{}}
\toprule
文献 & 吸收的核心思想 & 在实现中的落点 \\
\midrule
Vahrenkamp 等 2013
  & IRD = 固定 TCP，对基座位姿打分
  & \(\Tcp\) 固定（或缓变）；经 \(\Taxis(y_{\mathrm{rail}})\) 查询 \\
Rudorfer RM4D 2024
  & 商掉常见 yaw/roll 对称 $\to\sim$4D 图，RM 与 IRD 共用
  & flange 图卡 \(c\)；对 probe45 仅商掉基座 yaw \\
Murooka 等 2025
  & 可微可达图 \(f_{R}(r)\) 用于优化；
    分类器在可行域外梯度可能平坦
  & 单输出有符号 clearance；边界 stencil；查询全程 AD \\
Chiu 1988
  & 姿态应匹配任务的力/速度传递需求
  & Tip$\times$roll / 不确定锥把任务相容的自由或
    不确定集合编入 \(C\) \\
\bottomrule
\end{tabular}
\end{table}

\paragraph{Murooka 的可微可达图。}
Murooka 等定义标量 \(f_{R}(r)\)，满足
\begin{equation}
\label{eq:murooka}
  f_{R}(r)\ge 0 \text{ 若可达},\qquad
  f_{R}(r)< 0 \text{ 否则},
\end{equation}
并以光滑二分类器训练（例如 softplus logistic 损失）。
当 \(f_{R}\) 对任务坐标连续可微时，梯度 \(\partial f_{R}/\partial r\)
可直接进入轨迹优化而无需反复求解 IK。他们也指出：纯决策函数在远离可达集时
可能几乎为常数，从而给出无信息梯度。我们的设计沿用同一符号约定，但用局部
碰撞校验 stencil 与边界法向监督显式有符号 clearance，使 \(\nabla f\)
在边界附近仍有信息量（第~\ref{sec:training}~节）。

\paragraph{RM4D 对称性。}
Rudorfer 观察到：对许多 6/7 轴臂，可达性近似不变于绕重力的基座 yaw
以及绕末关节的 TCP roll，从而把 \(\SE(3)\) 查询降到固有的 4 维商空间。
我们沿用降维思路，但只强制 RM75 + probe45 真正成立的对称：基座（J1）yaw
是真自由对称；TCP/flange roll \emph{不是}，因为 TCP 的 \(z\) 轴相对 J7 轴
偏置约 \(\sim 50^{\circ}\)。

\paragraph{Chiu 的任务相容性。}
Chiu 度量姿态的速度/力椭球与任务的对齐程度。我们不在线计算 Yoshikawa 指标；
而是把任务自由度（大 tip$\times$roll 锥）与医学不确定（小盒$\times$锥）
编码为对 \(f\) 的集合查询，使 \(C\) 成为任务条件化的可达裕度。

% -----------------------------------------------------------------------------
\section{坐标系、导轨基座与 flange 图卡}
\label{sec:chart}

\subsection{世界位姿与相对位姿}
\label{sec:frames}

齐次位姿写作
\begin{equation}
\label{eq:pose-matrix}
  T
  \;=\;
  \begin{bmatrix}
    R & p \\
    0^{\mathsf T} & 1
  \end{bmatrix}
  \in\SE(3),
  \qquad
  R\in\SO(3),\; p\in\R^{3}.
\end{equation}
给定世界系 TCP \(\Tcp^{W}\) 与世界系 J1 轴坐标系 \(\Taxis^{W}\)，
TCP 在轴坐标系下表示为
\begin{equation}
\label{eq:relative-pose}
\begin{aligned}
  R_{\mathrm{axis\leftarrow tcp}}
  &=
  (R_{\mathrm{axis}}^{W})^{\mathsf T}\, R_{\mathrm{tcp}}^{W},
  \\
  p_{\mathrm{axis\leftarrow tcp}}
  &=
  (R_{\mathrm{axis}}^{W})^{\mathsf T}
  \bigl(p_{\mathrm{tcp}}^{W}-p_{\mathrm{axis}}^{W}\bigr).
\end{aligned}
\end{equation}
\emph{解释。}
式~\eqref{eq:relative-pose} 即群乘积 \((\Taxis^{W})^{-1}\Tcp^{W}\)。
所有可达查询都在该相对坐标系下提出，这正是 IRD 视角：移动基座/轴系时，
即使世界 TCP 固定，相对 TCP 也会改变。

\subsection{导轨诱导的轴坐标系}
\label{sec:rail}

对沿局部轴 \(e_{a}\) 的直线导轨（代码默认 \(a=1\)，即 \(y\)），
\begin{equation}
\label{eq:rail-base}
  \Taxis^{W}(y)
  \;=\;
  T_{\mathrm{world\leftarrow rail}}\;
  \underbrace{
    \begin{bmatrix}
      I_{3} & y\, e_{a} \\
      0^{\mathsf T} & 1
    \end{bmatrix}
  }_{\text{由导轨坐标 }y\text{ 产生的纯平移}}\;
  T_{\mathrm{rail\leftarrow axis}}(y{=}0).
\end{equation}
\emph{解释。}
\(y\) 是 URDF 原始导轨坐标。对~\eqref{eq:rail-base} 关于 \(y\) 求导
只沿导轨移动轴原点；旋转不变。这是把「相对基座放置的 IRD」
变成「相对单个导轨自由度的 IRD」的原语。

\subsection{TCP 到 flange}
\label{sec:tcp-flange}

设 \(T_{\mathrm{flange\leftarrow tcp}}\) 为 URDF 固定工具变换。
轴坐标系下 flange（link\_7）位姿为
\begin{equation}
\label{eq:tcp-to-flange}
  \bigl(p_{\mathrm{fl}},R_{\mathrm{fl}}\bigr)
  \;=\;
  \mathrm{compose}\bigl(
    (p_{\mathrm{tcp}},R_{\mathrm{tcp}}),\;
    T_{\mathrm{flange\leftarrow tcp}}
  \bigr).
\end{equation}
我们对 flange 几何而非 TCP 几何打分，使图卡坐标轴与 J7 转轴对齐
（下文不变量 / roll 分组需要）。

\subsection{九维 yaw 不变 flange 嵌入}
\label{sec:9d}

记 flange 位置 \(p=(p_{x},p_{y},p_{z})\) 与 flange 轴
\(u_{x},u_{y},u_{z}\)（\(R_{\mathrm{fl}}\) 的列）。定义径向光滑距离
\begin{equation}
\label{eq:radial}
  r
  \;=\;
  \sqrt{p_{x}^{2}+p_{y}^{2}+\varepsilon},
  \qquad
  \varepsilon=(10^{-3}\,\mathrm{m})^{2}.
\end{equation}
图卡映射为
\begin{equation}
\label{eq:chart}
  c(p,R)
  \;=\;
  \begin{pmatrix}
    p_{z}\\
    r\\
    u_{x}\!\cdot\!\hat z\\
    p_{xy}\!\cdot\! u_{x,xy}\\
    p_{xy}\!\times\! u_{x,xy}\\
    u_{z}\!\cdot\!\hat z\\
    p_{xy}\!\cdot\! u_{z,xy}\\
    p_{xy}\!\times\! u_{z,xy}\\
    u_{y}\!\cdot\!\hat z
  \end{pmatrix}
  \in\R^{9},
\end{equation}
其中 \(p_{xy}=(p_{x},p_{y})\)，\(u_{\cdot,xy}\) 为轴的 \(xy\) 投影，
\(\hat z=(0,0,1)\)，二维「叉积」为标量 \(p_{x}u_{y}-p_{y}u_{x}\)，且
\begin{equation}
\label{eq:uyz}
  u_{y}\!\cdot\!\hat z
  \;=\;
  (u_{z}\times u_{x})\!\cdot\!\hat z
  \;=\;
  u_{z,x}\,u_{x,y}-u_{z,y}\,u_{x,x}.
\end{equation}

\paragraph{各分量含义。}
\begin{itemize}[leftmargin=1.4em]
  \item \(p_{z}\)：沿 J1 轴的高度（对绕 \(\hat z\) 的 yaw 不变）。
  \item \(r\)：柱半径；\(\varepsilon\) 去掉轴上
        \(\sqrt{p_{x}^{2}+p_{y}^{2}}\) 的尖点
        （曲率约 \(10^{3}/\mathrm{m}\) 而非 \(10^{6}/\mathrm{m}\)）。
  \item 含 \(u_{z}\) 的分量：J7 不变块
        （代码下标 \(\{0,1,5,6,7\}\)），因绕 J7 旋转会转 \(u_{x},u_{y}\)
        但不转 \(u_{z}\)。
  \item 含 \(u_{x}\) 与 \(u_{y}\!\cdot\!\hat z\) 的分量：
        \emph{携带 roll} 的块。它们保留 flange roll \(\gamma\)，
        对 probe45 是真实自由度，不可商掉。
  \item 第九项 \(u_{y}\!\cdot\!\hat z\) 补全 \(R^{\mathsf T}\hat z\)。
        若无此项，当 \(r\to 0\) 时混合积消失，绕轴姿态重建留下二歧义。
\end{itemize}

\paragraph{对称性陈述。}
若 \(R_{z}(\alpha)\) 是同时作用于基座（等价地作用于相对 flange 位姿）的
绕 J1 轴旋转，则
\begin{equation}
\label{eq:yaw-invariance}
  c\bigl(R_{z}(\alpha)\,p,\; R_{z}(\alpha)\,R\bigr)
  \;=\;
  c(p,R).
\end{equation}
因此 \(c\) 把固有商 \(\SE(3)/\mathrm{Yaw}_{J1}\) 光滑嵌入 \(\R^{9}\)
（仅商掉 yaw 后固有 5 维 flange 图卡的冗余坐标）。

\begin{remark}[遗留的 5 维 TCP 图卡]
更早的嵌入
\(c_{\mathrm{TCP}}=[p_{z},\,u_{z}\!\cdot\!\hat z,\,\|p_{xy}\|,\,
p_{xy}\!\cdot\!u_{xy},\,p_{xy}\times u_{xy}]\)
同时商掉了 yaw 与 TCP roll。对 probe45 这不正确，仅保留用于审计。
\end{remark}

% -----------------------------------------------------------------------------
\section{有符号可达场}
\label{sec:field}

\subsection{定义}
\label{sec:field-def}

\begin{definition}[有符号 clearance 场]
网络 \(f_{\theta}:\R^{9}\to\R\) 在 flange 图卡上逼近有符号可达 clearance：
\begin{equation}
\label{eq:sign-convention}
  f_{\theta}(c)
  \;
  \begin{cases}
    >0 & \text{可达（正裕度）},\\
    =0 & \text{近似边界},\\
    <0 & \text{不可达}.
  \end{cases}
\end{equation}
\end{definition}
这与 Murooka 的符号模式~\eqref{eq:murooka} 一致，但\emph{幅度}被训练为跟踪
声明度量下经碰撞校验的局部距离（第~\ref{sec:metric}~节），而非仅分类间隔。

\subsection{输入归一化与傅里叶特征}
\label{sec:fourier}

由训练图卡行确定中心 \(\mu\in\R^{9}\) 与尺度 \(\sigma\in\R^{9}\)
（分位数中程）。归一化输入为
\begin{equation}
\label{eq:normalize}
  \tilde c
  \;=\;
  \frac{c-\mu}{\sigma}
  \qquad
  \text{（逐分量）}.
\end{equation}
\(B\) 个频带的傅里叶特征（默认 \(B{=}3\)）为
\begin{equation}
\label{eq:fourier}
  \Phi(\tilde c)
  \;=\;
  \Bigl[
    \tilde c,\;
    \sin(\pi 2^{0}\tilde c),\;
    \cos(\pi 2^{0}\tilde c),\;
    \ldots,\;
    \sin(\pi 2^{B-1}\tilde c),\;
    \cos(\pi 2^{B-1}\tilde c)
  \Bigr]
  \in\R^{9(1+2B)}.
\end{equation}
\emph{解释。}
米制通道与无量纲姿态通道动态范围不相容；若无~\eqref{eq:normalize}，
\(\sin(\pi 2^{k}x)\) 会打乱它们。傅里叶提升有助于表示尖锐的可达/不可达过渡，
而无需空间哈希网格（旧模型会记忆边界采样点）。

\subsection{网络结构}
\label{sec:mlp}

记 \(\mathrm{softplus}_{\beta}(z)=\beta^{-1}\log(1+e^{\beta z})\)。
场是带 Softplus 激活与线性标量头的残差 MLP：
\begin{equation}
\label{eq:mlp}
\begin{aligned}
  h_{0}
  &=
  \mathrm{softplus}_{\beta}\bigl(W_{0}\Phi(\tilde c)+b_{0}\bigr),
  \\
  h_{\ell+1}
  &=
  \mathrm{softplus}_{\beta}\bigl(
    h_{\ell}+W_{\ell}^{(2)}\,
    \mathrm{softplus}_{\beta}(W_{\ell}^{(1)}h_{\ell}+b_{\ell}^{(1)})
    +b_{\ell}^{(2)}
  \bigr),
  \\
  f_{\theta}(c)
  &=
  w^{\mathsf T}h_{L}+b.
\end{aligned}
\end{equation}
Softplus（不同于 ReLU）在有限输入上是 \(C^{\infty}\) 且导数处处非零，
当把 \(\nabla_{c}f\) 用作引导方向时更为理想。

\subsection{支撑域守卫}
\label{sec:guard}

设 \([\underline c,\overline c]\) 为训练图卡行的轴对齐包围盒。
归一化后，光滑外部距离为
\begin{equation}
\label{eq:support-dist}
  d_{\mathrm{supp}}(c)
  \;=\;
  \sqrt{
    \bigl\|
      \mathrm{ReLU}(\underline{\tilde c}-\tilde c)
      +
      \mathrm{ReLU}(\tilde c-\overline{\tilde c})
    \bigr\|_{2}^{2}
    +\epsilon^{2}
  }
  -\epsilon,
\end{equation}
守卫后的分数为
\begin{equation}
\label{eq:guarded}
  \tilde f(c)
  \;=\;
  f_{\theta}(c)
  -
  \lambda_{g}\, d_{\mathrm{supp}}(c),
  \qquad
  \lambda_{g}=8\text{（默认）}.
\end{equation}
\emph{解释。}
在训练支撑域外，\(\tilde f\) 被拉向负值，从而抑制优化外推到未标注图卡区域。

\subsection{世界系打分（IRD 点查询）}
\label{sec:score-world}

\begin{equation}
\label{eq:score-world}
  F(\Tcp^{W},\Taxis^{W})
  \;=\;
  \tilde f\Bigl(
    c\bigl(
      \mathrm{flange}\bigl(
        (\Taxis^{W})^{-1}\Tcp^{W}
      \bigr)
    \bigr)
  \Bigr).
\end{equation}
\emph{解释。}
固定 \(\Tcp^{W}\)，经~\eqref{eq:rail-base} 改变 \(\Taxis^{W}(y)\)：则
\(F\) 正是导轨/基座放置的光滑 IRD 分数。固定轴系、改变 \(\Tcp^{W}\)：
则 \(F\) 是光滑前向可达分数。相对图卡~\eqref{eq:chart} 使同一场服务两种视角。

% -----------------------------------------------------------------------------
\section{声明的 SE(3) 度量与训练目标}
\label{sec:training}

\subsection{度量}
\label{sec:metric}

边界标签使用将 \(1\,\mathrm{cm}\) 平移与 \(1^{\circ}\) 旋转等价的加权距离：
\begin{equation}
\label{eq:lambda-metric}
  \lambda
  \;=\;
  \frac{0.01\,\mathrm{m}}{1^{\circ}\text{（以弧度计）}}
  \;=\;
  \frac{0.01}{\pi/180}
  \;\approx\;
  0.5730\,\mathrm{m/rad}.
\end{equation}
对平移增量 \(\Delta p\in\R^{3}\)（米）与旋转增量 \(\Delta\omega\in\R^{3}\)
（弧度，旋转向量 / \(\mathfrak{so}(3)\) 坐标），
\begin{equation}
\label{eq:se3-distance}
  d(\Delta p,\Delta\omega)
  \;=\;
  \sqrt{
    \|\Delta p\|_{2}^{2}
    +
    \lambda^{2}\,\|\Delta\omega\|_{2}^{2}
  }.
\end{equation}
\emph{解释。}
这\emph{不是}双不变 \(\SE(3)\) 测地；它是声明的、图卡友好的度量，
用于放置 stencil 样本，并把 clearance 解释为「米」。

\subsection{Stencil 监督}
\label{sec:stencil}

在碰撞校验过的边界中心周围，于偏移
\begin{equation}
\label{eq:stencil-offsets}
  \pm 1,3,6,10\,\mathrm{mm}
  \quad\text{以及}\quad
  \pm 1,3,5^{\circ}
\end{equation}
处评估精确可达/不可达标签，并加入显式取零的二分边界中心。
它们为 \(f_{\theta}(c)\) 的有符号值提供目标 \(s^{\star}(c)\)。

\subsection{损失函数}
\label{sec:losses}

设小批量 \(\{(c_{i},y_{i},s_{i}^{\star},n_{i})\}_{i}\) 含二值标签
\(y_{i}\in\{0,1\}\)、SDF 目标 \(s_{i}^{\star}\)，以及（若有）归一化图卡坐标下的
单位边界法向 \(n_{i}\)。

\paragraph{分类（符号）。}
\begin{equation}
\label{eq:loss-cls}
  \mathcal{L}_{\mathrm{cls}}
  \;=\;
  \frac{
    \sum_{i} w_{i}^{\mathrm{cls}}\,
    \ell_{\mathrm{BCE}}\bigl(\alpha f_{\theta}(c_{i}),\, y_{i}\bigr)
  }{
    \sum_{i} w_{i}^{\mathrm{cls}}
  },
\end{equation}
其中 \(\alpha\) 为 logit 尺度（默认 \(3\)），
\(\ell_{\mathrm{BCE}}\) 为带 logits 的二元交叉熵。
\emph{解释。}
强制 Murooka 风格的符号正确性~\eqref{eq:sign-convention}。

\paragraph{有符号值（幅度）。}
对 SDF 权重 \(w_{i}^{\mathrm{sdf}}>0\) 的行，
\begin{equation}
\label{eq:loss-sdf}
  \mathcal{L}_{\mathrm{sdf}}
  \;=\;
  \mathrm{SmoothL1}_{\beta=0.1}
  \bigl(f_{\theta}(c_{i}),\, s_{i}^{\star}\bigr).
\end{equation}
\emph{解释。}
使 \(|f|\) 作为 clearance 有意义，而非任意分类间隔。

\paragraph{边界法向对齐。}
对需要梯度的 \(\tilde c\)，
\begin{equation}
\label{eq:loss-normal}
  \mathcal{L}_{\mathrm{n}}
  \;=\;
  \mathbb{E}_{i}\Bigl[
    1
    -
    \frac{
      \nabla_{\tilde c}f_{\theta}(c_{i})
      \;\cdot\;
      n_{i}
    }{
      \|\nabla_{\tilde c}f_{\theta}(c_{i})\|\,\|n_{i}\|
    }
  \Bigr].
\end{equation}
\emph{解释。}
余弦损失使场梯度与碰撞校验穿越方向对齐。这就是为何 \(\nabla F\)
在边界附近指向可达性增加，而非沿任意等值线。

\paragraph{可选 Eikonal（生产中关闭）。}
\begin{equation}
\label{eq:loss-eikonal}
  \mathcal{L}_{\mathrm{Eik}}
  \;=\;
  \mathrm{SmoothL1}\bigl(
    \|\nabla_{\tilde c}f\|,\;
    v^{\star}
  \bigr),
\end{equation}
其中 \(v^{\star}=1\) 或经验 stencil 斜率。工作空间归一化后的混合单位使
通用单位速度先验条件差；生产依赖~\eqref{eq:loss-sdf}--\eqref{eq:loss-normal}。

\paragraph{总和。}
\begin{equation}
\label{eq:loss-total}
  \mathcal{L}
  \;=\;
  \lambda_{\mathrm{cls}}\mathcal{L}_{\mathrm{cls}}
  +
  \lambda_{\mathrm{sdf}}\mathcal{L}_{\mathrm{sdf}}
  +
  \lambda_{\mathrm{n}}\mathcal{L}_{\mathrm{n}}
  +
  \lambda_{\mathrm{Eik}}\mathcal{L}_{\mathrm{Eik}}
  +
  \cdots
\end{equation}

\begin{remark}[为何不单用分类器？]
分类器可在 \(c\) 深入不可达区时仍保持高准确率却几乎为常数，于是
\(\nabla_{x}F\approx 0\)——恰在优化最需要恢复方向时（Murooka 的警告）。
有符号 stencil 值加法向对齐在边界带上保持可用斜率。
\end{remark}

% -----------------------------------------------------------------------------
\section{可微集合查询}
\label{sec:region}

所有集合查询共享同一模式：在名义 TCP 附近采样位姿，用~\eqref{eq:score-world}
打分，再以光滑聚合器归约。

\subsection{局部位姿扰动}
\label{sec:perturb}

在医学 TCP 坐标系（列向量 \([b,t,n]\)：副法向、切向、法向）中，
位置偏移 \(\delta p_{\mathrm{loc}}\in\R^{3}\) 与旋转向量
\(\delta\omega_{\mathrm{loc}}\in\R^{3}\) 产生
\begin{equation}
\label{eq:perturb}
\begin{aligned}
  p'
  &=
  p + R\,\delta p_{\mathrm{loc}},
  \\
  R'
  &=
  R\;\operatorname{Exp}(\delta\omega_{\mathrm{loc}}),
\end{aligned}
\end{equation}
其中 \(\operatorname{Exp}:\R^{3}\to\SO(3)\) 为 Rodrigues 映射
\begin{equation}
\label{eq:rodrigues}
  \operatorname{Exp}(\omega)
  \;=\;
  I
  +
  \frac{\sin\theta}{\theta}[\hat\omega]
  +
  \frac{1-\cos\theta}{\theta^{2}}[\hat\omega]^{2},
  \qquad
  \theta=\|\omega\|,
\end{equation}
\([\hat\omega]\) 为单位轴的反对称矩阵
（\(\theta\to 0\) 时极限为 \(I+[\omega]\)）。

\subsection{归一化 softmin 与 softmax}
\label{sec:soft}

对向量 \(v\in\R^{K}\) 与温度 \(\tau>0\)，
\begin{equation}
\label{eq:logsumexp}
  \LogSumExp(v)
  \;=\;
  \log\sum_{k=1}^{K} e^{v_{k}}.
\end{equation}
用于鲁棒聚合的\emph{归一化 softmin} 为
\begin{equation}
\label{eq:softmin}
  \softmin_{\tau}(v)
  \;=\;
  -\tau\Bigl(
    \LogSumExp(-v/\tau)
    -
    \log K
  \Bigr).
\end{equation}
\emph{为何要有 \(\log K\) 项？}
若对所有 \(k\) 有 \(v_{k}\equiv a\)，则
\(\LogSumExp(-v/\tau)=\log K - a/\tau\)，故 \(\softmin_{\tau}(v)=a\)。
常数场保持不变；若无 \(\log K\)，softmin 会偏移 \(-\tau\log K\)。

\emph{归一化 softmax 型最大值}为
\begin{equation}
\label{eq:softmax}
  \softmax_{\tau}(v)
  \;=\;
  \tau\Bigl(
    \LogSumExp(v/\tau)
    -
    \log K
  \Bigr).
\end{equation}
对最大值同样具有常数保持性质。

\paragraph{隐式权重（梯度结构）。}
对~\eqref{eq:softmin} 求导得
\begin{equation}
\label{eq:softmin-grad}
  \frac{\partial\softmin_{\tau}(v)}{\partial v_{k}}
  \;=\;
  \frac{e^{-v_{k}/\tau}}{\sum_{j}e^{-v_{j}/\tau}}
  \;=\;
  \pi_{k}^{\mathrm{min}}(v),
\end{equation}
即集中在\emph{最小}分量上的 Gibbs 分布。对~\eqref{eq:softmax}，
\begin{equation}
\label{eq:softmax-grad}
  \frac{\partial\softmax_{\tau}(v)}{\partial v_{k}}
  \;=\;
  \frac{e^{v_{k}/\tau}}{\sum_{j}e^{v_{j}/\tau}}
  \;=\;
  \pi_{k}^{\mathrm{max}}(v),
\end{equation}
集中在\emph{最大}分量。因此鲁棒查询把梯度回传到最差样本；
任务自由查询把梯度回传到最佳样本。

\subsection{\texorpdfstring{Region A：不确定盒 $\times$ 小锥（softmin）}{Region A：不确定盒 x 小锥（softmin）}}
\label{sec:region-a}

\begin{definition}[Region A]
默认范围：副法向 \(4\,\mathrm{mm}\)、切向 \(3\,\mathrm{mm}\)、法向 \(2\,\mathrm{mm}\)；
姿态锥半角 \(\beta=3^{\circ}\)；
\(K=64\) 个固定 scrambled Sobol 样本（含名义位姿）。
\end{definition}

位置半样本使用联合 5 维 Sobol 序列 \(u\in(0,1)^{5}\)：
\begin{equation}
\label{eq:region-pos}
  \delta p
  \;=\;
  (2u_{1:3}-1)\odot
  (b,t,n)_{\mathrm{extent}},
\end{equation}
姿态样本在锥内对立体角均匀，
\begin{equation}
\label{eq:region-ori}
\begin{aligned}
  \cos\rho
  &=
  1-u_{4}\bigl(1-\cos\beta\bigr),
  \\
  \rho
  &=
  \arccos(\cos\rho),
  \\
  \phi
  &=
  2\pi u_{5},
  \\
  \delta\omega
  &=
  (\rho\cos\phi,\;\rho\sin\phi,\;0),
\end{aligned}
\end{equation}
再反对称化为 \(\{\,0,\,\pm\text{半样本}\,\}\) 至长度 \(K\)。
Region~A clearance 为
\begin{equation}
\label{eq:region-a}
  C_{\mathrm{A}}(\Tcp,\Taxis)
  \;=\;
  \softmin_{\tau}
  \Bigl\{
    F\bigl(\Tcp(\delta_{k}),\,\Taxis\bigr)
  \Bigr\}_{k=1}^{K},
  \qquad
  \tau=0.15.
\end{equation}
\emph{解释。}
\(C_{\mathrm{A}}\) 是医学动机不确定集 \(U\) 上的光滑下包络。
提高 \(C_{\mathrm{A}}\) 即改善小配准/接触误差下最坏可能位姿。
覆盖率诊断使用 \(\frac{1}{K}\sum_{k}\sigma(F_{k}/\tau)\)。

\subsection{\texorpdfstring{任务锥：tip $\times$ roll 自由集（softmax）}{任务锥：tip x roll 自由集（softmax）}}
\label{sec:task-cone}

超声接触常允许较大 tip 倾斜与探头 roll，它们是\emph{控制器可选}的自由变量，
而非不确定量。

\begin{definition}[TaskCone]
Tip 半角 \(\beta_{\mathrm{tip}}\)（例如 \(20^{\circ}\)--\(45^{\circ}\)），
roll 半程 \(\gamma_{\max}\)（例如 \(20^{\circ}\)--\(30^{\circ}\)），
在 \((\rho,\phi,\gamma)\) 上取 \(K\) 个 Sobol 样本。
\end{definition}
\begin{equation}
\label{eq:task-cone-sample}
\begin{aligned}
  \cos\rho
  &=
  1-u_{1}\bigl(1-\cos\beta_{\mathrm{tip}}\bigr),
  \\
  \phi
  &=
  2\pi u_{2},
  \\
  \gamma
  &=
  (2u_{3}-1)\gamma_{\max},
  \\
  \delta\omega
  &=
  (\rho\cos\phi,\;\rho\sin\phi,\;\gamma).
\end{aligned}
\end{equation}
位置固定（接触点给定）。任务 clearance 为
\begin{equation}
\label{eq:task-cone}
  C_{\mathrm{task}}(\Tcp,\Taxis)
  \;=\;
  \softmax_{\tau}
  \Bigl\{
    F\bigl(\Tcp(\eta_{k}),\,\Taxis\bigr)
  \Bigr\}_{k=1}^{K},
  \qquad
  \tau=0.20.
\end{equation}
\emph{解释。}
与 Region~A 相反：\(C_{\mathrm{task}}\) 询问是否\emph{存在} tip$\times$roll
选择使手臂舒适。梯度把基座/导轨推向锥内存在高 clearance 自由位姿的放置。
这是 Chiu 式「让姿态集合匹配任务」的思想，但用可达裕度而非椭球指标表达。

\subsection{\texorpdfstring{嵌套自由 $\times$ 不确定查询（SetQuery）}{嵌套自由 x 不确定查询（SetQuery）}}
\label{sec:set-query}

Phase-4 一般语义为
\begin{equation}
\label{eq:nested-semantic}
  C_{\mathrm{set}}(T)
  \;=\;
  \max_{\eta\in D_{\mathrm{free}}}
  \;
  \underbrace{
    \CVaR_{\alpha}\text{ 或 }\softmin
  }_{\delta\in U}
  \;
  F\bigl(c(T(\eta,\delta))\bigr).
\end{equation}
批量分数形状为 \([N,K_{\mathrm{free}},K_{\mathrm{unc}}]\)。
先归约不确定轴，再对自由变量取光滑最大值。

\paragraph{Rockafellar--Uryasev 下尾 CVaR。}
对 clearance 值（越大越好），水平 \(\alpha\) 的下尾 CVaR 为
\begin{equation}
\label{eq:cvar}
  \CVaR_{\alpha}(X)
  \;=\;
  \min_{t\in\R}
  \Biggl\{
    t
    +
    \frac{1}{1-\alpha}\,
    \mathbb{E}\bigl[(t-X)_{+}\bigr]
  \Biggr\}.
\end{equation}
\emph{解释。}
\(t\) 扮演 VaR 阈值；期望惩罚低于 \(t\) 的短缺。
经验上最优 \(t\) 落在样本集中，故代码对每个样本值评估目标，再对这些候选
\(t\) 取硬或软最小值。当 \(K_{\mathrm{unc}}\) 中等时，这比纯 softmin 更统计稳定。

\paragraph{实现中的归约。}
\begin{equation}
\label{eq:set-query}
\begin{aligned}
  g(\eta)
  &=
  \CVaR_{\alpha}\bigl\{F(T(\eta,\delta_{j}))\bigr\}_{j}
  \quad\text{或}\quad
  \softmin_{\tau}\{\ldots\},
  \\
  C_{\mathrm{set}}
  &=
  \softmax_{\tau}\bigl\{g(\eta_{i})\bigr\}_{i}.
\end{aligned}
\end{equation}

% -----------------------------------------------------------------------------
\section{为何 IRD 梯度存在且有用}
\label{sec:gradients}

\subsection{IRD 点查询上的链式法则}
\label{sec:chain}

考虑世界 TCP 固定时的导轨坐标 \(y\)。则
\begin{equation}
\label{eq:chain-rail}
  \frac{\mathrm{d}F}{\mathrm{d}y}
  \;=\;
  \frac{\partial\tilde f}{\partial c}\;
  \frac{\partial c}{\partial (p,R)}\;
  \frac{\partial (p,R)}{\partial \Taxis}\;
  \frac{\mathrm{d}\Taxis}{\mathrm{d}y}.
\end{equation}
每一因子均由 Torch 算子实现（\(\operatorname{Exp}\)、矩阵乘、
\(\sqrt{\,\cdot+\varepsilon}\)、Softplus MLP），故 \(\mathrm{d}F/\mathrm{d}y\)
在浮点误差内精确。导轨与 TCP-\(x\) 上 AD 相对 FD 的中位相对误差约为 \(10^{-4}\)。

\paragraph{几何解读。}
\(\partial c/\partial(p,R)\) 把相对 flange 位姿的空间 twist 转为图卡坐标。
\(\partial\tilde f/\partial c\) 被训练为在边界附近与外指可达法向对齐
（式~\eqref{eq:loss-normal}）。因此 \(\mathrm{d}F/\mathrm{d}y>0\) 表示
「滑台沿 \(+y\) 移动会提高预测 clearance」。

\subsection{条件锥查询的梯度}
\label{sec:cond-grad}

对 Region~A，
\begin{equation}
\label{eq:grad-region}
  \nabla_{x} C_{\mathrm{A}}
  \;=\;
  \sum_{k=1}^{K}
  \pi_{k}^{\mathrm{min}}\,
  \nabla_{x} F\bigl(\Tcp(\delta_{k};x),\,\Taxis(x)\bigr).
\end{equation}
对 TaskCone，
\begin{equation}
\label{eq:grad-task}
  \nabla_{x} C_{\mathrm{task}}
  \;=\;
  \sum_{k=1}^{K}
  \pi_{k}^{\mathrm{max}}\,
  \nabla_{x} F\bigl(\Tcp(\eta_{k};x),\,\Taxis(x)\bigr).
\end{equation}
\emph{「条件 IRD 梯度」的解释。}
条件即集合 \(U\) 或 \(D_{\mathrm{free}}\)。标量 \(C\) 是集合聚合后的 clearance。
其对基座/导轨/路径参数的梯度，是采样局部位姿处\emph{点 IRD 梯度的加权平均}，
权重为~\eqref{eq:softmin-grad} 或~\eqref{eq:softmax-grad}。
无需离散体素 \(\arg\max\)。

\subsection{\texorpdfstring{血管表面例子 $(\theta,y)$}{血管表面例子 (theta, y)}}
\label{sec:vessel-grad}

在椭圆血管上，TCP 位置由表面角 \(\theta\) 与路径横坐标 \(s\) 约束，
导轨 \(y\) 设定基座。Demo 代码构造
\begin{equation}
\label{eq:demo-query}
  C(s;\theta,y)
  \;=\;
  C_{\mathrm{task}}\bigl(
    \Tcp(\theta,s),\,
    \Taxis(y)
  \bigr)
\end{equation}
并用 \(\mathtt{torch.autograd.grad}\) 提取
\begin{equation}
\label{eq:demo-grads}
  \frac{\partial C}{\partial\theta},
  \qquad
  \frac{\partial C}{\partial y}.
\end{equation}
这些分量的 quiver 图即 \((\theta,y)\) 规划平面上条件 IRD 向量场的直观可视化。

% -----------------------------------------------------------------------------
\section{轨迹算子与在线更新}
\label{sec:traj}

\subsection{聚合顺序}
\label{sec:agg-order}

给定路点 \(\{\Tcp^{(\ell)}\}_{\ell=1}^{W}\) 与（共享或逐点）轴坐标系，
轨迹算子按以下顺序作用：
\begin{enumerate}[leftmargin=1.6em]
  \item 对每个近似角度候选做\textbf{不确定 / Region~A 场景}；
  \item \textbf{角度聚合}（\(\softmax\) 选择、\(\softmin\) 鲁棒、均值或 CVaR）；
  \item \textbf{路径残差}：每段上的内部 \(\SE(3)\) 采样；
  \item \textbf{轨迹归约}：对路点与路径裕度做 softmin / CVaR / 硬最小。
\end{enumerate}

\subsection{SE(3) 路径插值}
\label{sec:se3-interp}

在从 \(T_{0}\) 到 \(T_{1}\) 的段上，对 \(\alpha\in(0,1)\)，
\begin{equation}
\label{eq:se3-interp}
\begin{aligned}
  p(\alpha)
  &=
  (1-\alpha)p_{0}+\alpha p_{1},
  \\
  R(\alpha)
  &=
  R_{0}\,
  \operatorname{Exp}\bigl(
    \alpha\,\operatorname{Log}(R_{0}^{\mathsf T}R_{1})
  \bigr).
\end{aligned}
\end{equation}
\emph{解释。}
平移线性运动；姿态沿 \(\SO(3)\) 测地。要求整段插值路径均可 IK 求解的控制器
需要这些内部样本：仅端点可达不足够。

\subsection{逐路点与路径裕度}
\label{sec:path-margin}

设 \(A\) 索引角度候选，\(S\) 索引 Region~A 场景。打分后，
\begin{equation}
\label{eq:waypoint-clearance}
\begin{aligned}
  g_{\ell,a}
  &=
  \mathcal{A}_{S}\bigl[F(T^{(\ell)}_{a,s})\bigr],
  \\
  C_{\ell}
  &=
  \mathcal{A}_{A}\bigl[g_{\ell,a}\bigr],
\end{aligned}
\end{equation}
内部路径样本 \(C^{\mathrm{path}}_{m}\) 类似。标量轨迹 clearance 为
\begin{equation}
\label{eq:traj-clearance}
  C_{\mathrm{traj}}
  \;=\;
  \softmin_{\tau}
  \Bigl[
    \{C_{\ell}\}_{\ell=1}^{W}
    \;\cup\;
    \{C^{\mathrm{path}}_{m}\}_{m}
  \Bigr]
\end{equation}
（或 CVaR / 硬最小）。可选的 SRS 臂角 \(\psi\) 在暴露时作为外层
\(\softmax\) 进入 \(\psi\) 条件场。

\subsection{为何景观随轨迹更新}
\label{sec:why-update}

网络 \(f_{\theta}\) 在 flange 图卡上训练一次后冻结。
轨迹只是曲线 \(x\mapsto(\Tcp(x),\Taxis(x))\)。由于：
\begin{enumerate}[leftmargin=1.6em]
  \item \(c\) 依赖\emph{相对}位姿~\eqref{eq:relative-pose}；
  \item 锥样本表达在\emph{当前} TCP 体坐标系~\eqref{eq:perturb} 中；
  \item 聚合器~\eqref{eq:softmin}--\eqref{eq:softmax} 是当前样本分数的无记忆函数；
\end{enumerate}
在新路径上重新评估 \(C_{\mathrm{task}}(s)\) 或 \(C_{\mathrm{traj}}(x)\)
会自动得到新的 IRD 景观与新梯度 \(\nabla_{x}C\)，\emph{无需}修改 \(\theta\)。

\subsection{典型优化目标}
\label{sec:opt}

椭圆血管 demo 最小化朝向共形 clearance 目标 \(m^{\star}\) 的铰链项加正则：
\begin{equation}
\label{eq:opt-loss}
\begin{aligned}
  \ell_{\mathrm{IRD}}
  &=
  \frac{1}{W}\sum_{\ell}
  \mathrm{softplus}\!\Bigl(\frac{m^{\star}-C_{\ell}}{\sigma}\Bigr)
  +
  \max_{\ell}\mathrm{softplus}\!\Bigl(\frac{m^{\star}-C_{\ell}}{\sigma}\Bigr)
  -
  \kappa\,\overline{C},
  \\
  \ell
  &=
  w_{\mathrm{IRD}}\,\ell_{\mathrm{IRD}}
  +
  w_{\mathrm{near}}\,\ell_{\mathrm{nearest}}
  +
  w_{\mathrm{cont}}\,\ell_{\mathrm{continuity}}
  +
  \cdots
\end{aligned}
\end{equation}
\emph{解释。}
\(\mathrm{softplus}((m^{\star}-C)/\sigma)\) 是光滑单边惩罚：
当 \(C\ge m^{\star}\) 时不激活，clearance 低于目标时上升。
额外的 \(-\kappa\overline{C}\) 项在铰链已满足后仍奖励向 IRD 峰攀升。
连续性惩罚保持 \(\theta(s)\) 与 \(y(s)\) 光滑。对~\eqref{eq:opt-loss}
做梯度下降，正是「用条件 IRD 梯度更新轨迹」。

% -----------------------------------------------------------------------------
\section{端到端交换图}
\label{sec:diagram}

\[
\begin{array}{c}
x
\;\xrightarrow{\;\text{运动学 / 导轨}\;}
(\Tcp(x),\Taxis(x))
\;\xrightarrow{\;\text{局部采样}\;}
\{T_{k}(x)\}
\\[0.6em]
\xrightarrow{\;\eqref{eq:relative-pose}+\eqref{eq:chart}\;}
\{c_{k}\}
\;\xrightarrow{\;\tilde f\;}
\{F_{k}\}
\;\xrightarrow{\;\mathcal{A}\;}
C(x)
\;\xrightarrow{\;\nabla\;}
\nabla_{x}C
\end{array}
\]

\begin{proposition}[单一自动微分磁带]
若上图每一箭头均由初等可微映射实现，则 \(\nabla_{x}C\) 几乎处处存在，
并与解析链式法则一致。离散 IRD 体素查表在表索引处打断磁带；神经场不会。
\end{proposition}

% -----------------------------------------------------------------------------
\section{范围、保证与非目标}
\label{sec:scope}

\begin{itemize}[leftmargin=1.4em]
  \item \textbf{引导，非证书。}
        留出边界误差在括号化时可达亚毫米，但该场是优化方向算子。
        硬可行性仍需最终多种子 IK 与碰撞检查。
  \item \textbf{对称近似。}
        仅 J1 yaw 精确；残余探头 roll 碰撞差异少见，由下游拒绝。
  \item \textbf{点场中无分支 / 连续 \(\psi\)。}
        离线标签是点可达性。路径残差与连续臂角选择在轨迹算子中处理。
  \item \textbf{本身不是完整患者轨迹规划器。}
        Region~A / TaskCone / TrajectoryTask 是点场上的查询层；
        应用 demo 另加几何路径参数化与损失项。
\end{itemize}

% -----------------------------------------------------------------------------
\section{符号表}
\label{sec:symbols}

\begin{center}
\begin{tabular}{@{}ll@{}}
\toprule
符号 & 含义 \\
\midrule
\(\Tcp,\Taxis\) & 世界 TCP 位姿；世界 J1 轴位姿 \\
\(c\in\R^{9}\) & yaw 不变 flange 图卡 \\
\(f_{\theta},\tilde f\) & 有符号场；支撑域守卫后的场 \\
\(F\) & 世界点分数~\eqref{eq:score-world} \\
\(C_{\mathrm{A}},C_{\mathrm{task}},C_{\mathrm{set}},C_{\mathrm{traj}}\) &
  聚合后的 clearance \\
\(y\) & 导轨坐标 \\
\(\delta,\eta\) & 不确定样本；自由（任务）样本 \\
\(\tau\) & softmin/softmax 温度 \\
\(\lambda\) & 平移--旋转度量尺度~\eqref{eq:lambda-metric} \\
\(\pi^{\mathrm{min}},\pi^{\mathrm{max}}\) & softmin / softmax 权重 \\
\bottomrule
\end{tabular}
\end{center}

% -----------------------------------------------------------------------------
\section{参考文献}
\label{sec:refs}

\begin{enumerate}[leftmargin=1.6em]
  \item N.~Vahrenkamp, T.~Asfour, and R.~Dillmann,
        ``Robot placement based on reachability inversion,'' ICRA 2013.
  \item M.~Rudorfer,
        ``RM4D: A Combined Reachability and Inverse Reachability Map for
        Common 6-/7-axis Robot Arms by Dimensionality Reduction to 4D,''
        ICRA 2025 / arXiv:2410.06968.
  \item M.~Murooka et~al.,
        ``Learning Differentiable Reachability Maps for Optimization-based
        Humanoid Motion Generation,'' Humanoids 2025 / arXiv:2508.11275.
  \item S.~L.~Chiu,
        ``Task Compatibility of Manipulator Postures,''
        International Journal of Robotics Research, 1988.
  \item 内部设计说明：
        \texttt{ird\_playground/docs/rm4d\_signed\_ird\_design.md}。
  \item 英文对照版：
        \texttt{ird\_playground/docs/ird\_operator\_mathematics.tex}。
\end{enumerate}

\appendix
\section{代码与公式对照}
\label{sec:code-map}

\begin{center}
\begin{tabular}{@{}ll@{}}
\toprule
公式 & 主要模块 \\
\midrule
\eqref{eq:rail-base}
  & \texttt{region/operator.py::base\_from\_rail\_torch} \\
\eqref{eq:chart}
  & \texttt{ird/canonical.py::canonical\_flange\_invariants\_torch} \\
\eqref{eq:score-world}
  & \texttt{neural/signed\_field.py::score\_world} \\
\eqref{eq:softmin}--\eqref{eq:softmax}
  & \texttt{region/operator.py}, \texttt{region/set\_query.py} \\
\eqref{eq:region-a}
  & \texttt{region/operator.py::RegionA} \\
\eqref{eq:task-cone}
  & \texttt{region/task\_cone.py::TaskConeReachability} \\
\eqref{eq:nested-semantic}--\eqref{eq:set-query}
  & \texttt{region/set\_query.py::SetQueryOperator} \\
\eqref{eq:se3-interp}--\eqref{eq:traj-clearance}
  & \texttt{region/trajectory\_operator.py} \\
\eqref{eq:loss-cls}--\eqref{eq:loss-total}
  & \texttt{neural/train\_signed.py} \\
\eqref{eq:opt-loss}
  & \texttt{experiments/ellipse\_vessel\_ird\_demo.py} \\
\bottomrule
\end{tabular}
\end{center}

\end{document}
